% ---
% RESUMOS
% ---

% RESUMO em português
\setlength{\absparsep}{18pt} % ajusta o espaçamento dos parágrafos do resumo
\begin{resumo}
	O projeto \textit{Civitas} consiste em um sistema web desenvolvido para apoiar a gestão pública municipal, com foco no controle administrativo e financeiro de órgãos públicos. A proposta surge diante da necessidade de centralizar informações, reduzir inconsistências cadastrais, melhorar o acompanhamento de despesas e ampliar a rastreabilidade das operações relacionadas ao orçamento público. O sistema oferece uma plataforma digital para organização de cadastros administrativos, registro e acompanhamento de orçamentos, despesas, documentos, fluxos e auditorias, além de recursos de consulta, atualização e controle da situação dos registros. A solução também disponibiliza um painel com indicadores financeiros, permitindo melhor visualização das informações e apoio à tomada de decisão pelos gestores. Em seus principais aspectos técnicos, o \textit{Civitas} utiliza tecnologias voltadas ao desenvolvimento web, integrando interface de usuário, serviços de aplicação, banco de dados e mecanismos de autenticação, de forma a proporcionar segurança, organização e possibilidade de evolução do sistema. Como resultado, espera-se que o projeto contribua para maior agilidade nos processos administrativos, redução de erros, melhoria do controle interno, fortalecimento da transparência e gestão mais eficiente dos recursos públicos municipais.
	
	\noindent\textbf{Palavras-chave}: gestão pública municipal; controle orçamentário; despesas públicas; transparência administrativa.
\end{resumo}

% ABSTRACT in english
\begin{resumo}[Abstract]
	\begin{otherlanguage*}{english}
		The \textit{Civitas} project consists of a web-based system developed to support municipal public management, focusing on administrative and financial control of public agencies. The proposal arises from the need to centralize information, reduce data inconsistencies, improve expense tracking, and enhance the traceability of operations related to the public budget. The system provides a digital platform for organizing administrative records, registering and monitoring budgets, expenses, documents, workflows, and audits, as well as offering features for querying, updating, and controlling the status of records. The solution also includes a dashboard with financial indicators, enabling better visualization of information and supporting decision-making by managers. In its main technical aspects, the \textit{Civitas} uses technologies aimed at web development, integrating user interface, application services, database, and authentication mechanisms, in order to provide security, organization, and scalability for the system. As a result, the project is expected to contribute to greater agility in administrative processes, reduction of errors, improvement of internal control, strengthening of transparency, and more efficient management of municipal public resources.
		
		\noindent\textbf{Keywords}: municipal public management; budget control; public expenses; administrative transparency.
	\end{otherlanguage*}
\end{resumo}